\section{The main-scale saturation domain}\label{sec:main-saturation}
\label{sat:section}

Six zeros force one derivative crossing before the events, one between
them and three after them. Their alternating values constrain the controls
and states. We work in each fixed permitted family, without a rank
assumption and with common reductions independent of relative rates.
First record the tail estimate that will complete weighted confinement.

\begin{lemma}[A tail bound from weighted controls]\label{mainloc:tail}
For every fixed $K_0>0$, there are common constants $M,m_0>0$ such
that, whenever $0<|m|<m_0$ and
\[
 |a|\le K_0m^4,\qquad |b|\le K_0m^2,\qquad
 |\delta|\le K_0|m|^3,
\]
every zero of the actual displacement in the common physical interval
$I$ satisfies $|q|<Mm^2$. On the two complementary tails, $D$ has
the sign of $c$. No observation-rank or root-count hypothesis is needed.
\end{lemma}

\begin{proof}
Write $v=Q-q$. For $q\ge Mm^2$, reduce $m_0$ so that
$|\delta|/q\le K_0|m|/M\le1$. Equation~\eqref{eq:quartic-global-value}
gives
\[
 \frac{|v|}{q^2}\le K\{M^{-1/2}+M^{-3/2}+m^2+\sqrt{q_0}
                                      +m^2M^{-1/2}\}.
\]
In particular, the mismatch and separation terms use respectively
$|m|/\sqrt q\le M^{-1/2}$ and
$|\delta|\sqrt{q+|\delta|}/q^2\le KM^{-3/2}$.
For small $q_0$, this also gives $|v|/q<1/2$. The exact closing
expansion yields
\[
 D-cq^2=a+bq+(1+b)v+c(2qv+v^2)+Q^3S(Q,\lambda).
\]
Consequently
\begin{equation}\label{sat:tail-estimate}
 \frac{|D(q)-c(\lambda)q^2|}{q^2}
 \le K\{M^{-2}+M^{-1}+M^{-1/2}+M^{-3/2}
                              +\sqrt{q_0}+|m|\}.
\end{equation}
The first two terms bound $a/q^2$ and $b/q$; the remaining closing
terms follow from $|Q|\le2q$ and the bound for $v$.
For $q\le-Mm^2$, both passage inputs are negative once $m_0$ is
small: $q<\min(0,\delta)$. Thus $Q=q$ exactly, and the same estimate
follows directly from the closing expansion. Choose $M$ large, then
$q_0$ small and $m_0$ small so that the right-hand side is less than
half the common lower bound for $|c|$.

These estimates can be made on a smaller inner interval without losing
the original common $I$. At the reference, \eqref{eq:center-displacement}
gives $D(0)=D'(0)=0$ and strict curvature of sign $c_0$ throughout
the already reduced common interval. Hence $D$ has that sign at every
nonzero reference state. A closed outer annulus has a nonzero value
margin; joint continuity in the parameters preserves it after one common
parameter reduction. The common domains are chosen to contain
$\overline I$, so this argument covers the whole original interval.
All choices precede the choice of a control path.
\end{proof}

\begin{lemma}[Saturation of the curvature count]
\label{sat:saturation}\label{corner:saturated-geometry}
\label{cor:six-necessary-window}
If the actual displacement has six distinct zeros in the common interval,
then
\[
 mc_0(p)<0,\qquad \delta A(p)c_0(p)<0,\qquad
 0<\delta/\delta_{\rm c}(m,p)<1.
\]
In particular $0<|\delta|<|\delta_{\rm c}|\le K|m|^3$ uniformly.
Put $s=-\sgn c_0$, $E=sD$, and denote the events by $e_1<e_2$.
There are two simple outer curvature zeros $\zeta_1<\zeta_2$, and
\begin{equation}\label{sat:curvature-signs}
 \begin{array}{c|ccccc}
 q&(-q_0,e_1)&(e_1,e_2)&(e_2,\zeta_1)&
       (\zeta_1,\zeta_2)&(\zeta_2,q_0)\\ \hline
 \sgn E''&-&+&-&+&-
 \end{array}
\end{equation}
The continuous $E'$ has exactly one zero in each open interval, none
at its shared junctions, and
\begin{equation}\label{sat:junction-slopes}
 E'(e_1)<0,\quad E'(e_2)>0,\quad
 E'(\zeta_1)<0,\quad E'(\zeta_2)>0.
\end{equation}
All six zeros are simple $C^1$ zeros, including any at a threshold.
The extremal values have signs $+,-,+,-,+$; the first is an inactive
positive maximum and the second a negative minimum between the events.
The window is necessary, not a sufficient chamber. On
$\delta=\delta_{\rm c}$, $mc_0<0$, the distinct-point bound is four.
\end{lemma}

\begin{proof}
The four-point slices exclude $m=0$ and $\delta=0$.
The event signs~\eqref{eq:event-signs} give
$s_E=1+\mathbf1_{\{\delta Ac_0<0\}}$ threshold curvature changes;
there are at most $s_O=2$ outer changes. Joining same-sign pieces
through isolated zeros and events gives at most $s_E+s_O+1\le5$
strict monotonicity blocks for the continuous $D'$. Each block,
with its finite endpoints, contains at most one zero. Six zeros of
$D$ give five distinct derivative zeros in their open gaps by $C^1$
Rolle, so $s_E=s_O=2$. Proposition~\ref{prop:cubic-curvature-window}
gives the stated window and cubic bound. The event signs and outer
factorization~\eqref{eq:curvature-factor} give~\eqref{sat:curvature-signs}.

Retain every shared finite junction in both adjacent blocks of $E'$.
A derivative zero there would occupy both blocks, leaving room for at
most four distinct derivative zeros. Hence all five lie in the five
respective interiors, and strict monotonicity gives
\eqref{sat:junction-slopes}. A stationary fixed point would supply a
sixth derivative zero distinct from the five open-gap zeros. Thus none
occurs, including at an event, without assigning it a smooth multiplicity.
The extrema interlace the six roots, giving all five strict value signs.
At $\delta=\delta_{\rm c}$ the single double outer curvature zero does
not change sign, leaving at most three blocks and four distinct fixed
points. Truncating at the ends of the common interval removes blocks
but cannot add them.
\end{proof}

\begin{proposition}[Weighted confinement]
\label{sat:weighted-confinement}
For every fixed family satisfying the hypotheses of
Theorem~\ref{thm:six-main}, there are a further common reduction of
the actual return interval and parameter neighbourhood and a constant
\(K>0\) with the following property. If a member has six distinct
fixed points \(q_1,\ldots,q_6\), then
\begin{equation}\label{sat:weighted-bounds}
 0<|\delta|\le K|m|^3,\qquad
 |b|\le K|m|^2,\qquad |a|\le K|m|^4,\qquad
 \max_j|q_j|\le K|m|^2.
\end{equation}
The interval, parameter bounds and constant are independent of the
relative rates of the observations. No rank condition on those
observations is required. In the rank-free recognition formulation,
the same bounds hold with the canonical observed strength
\(\widehat m=C_{1,1}(\lambda)+C_{1,2}(\lambda)\) in place of
the independent chart variable \(m\).
\end{proposition}

\begin{proof}
Use Lemma~\ref{lem:auxiliary-controls} when strength and separation
are dependent observations, and restrict back to the original family
at \(m=\widehat m\) at the end. All estimates concern smooth root profiles on common
compact reductions. The regular jets \(a,b\) are always the exact
observations \eqref{eq:regular-jets}.

Lemma~\ref{cor:six-necessary-window} already gives the first
bound in \eqref{sat:weighted-bounds}. At a curvature zero write
\(q=r^2\). The exact division \eqref{eq:V-division} gives
\begin{equation}\label{sat:curvature-state}
 \delta=r^2\{mH(r,m,p)+rB(r,p)\}.
\end{equation}
Here \(H\) has sign \(A\), and \(B\) has sign \(c_0/A\),
with uniform nonzero margins. In the sector forced by six zeros,
\(mH\) has the sign of \(\delta\), whereas \(rB\) has the
opposite sign. Equation~\eqref{sat:curvature-state} consequently
implies
\begin{equation}\label{sat:curvature-scale}
 r\le K|m|,\qquad
 |v|=|\delta|/r^2=|mH+rB|\le K|m|
\end{equation}
at both curvature zeros. This argument does not require a lower
bound for either positive root \(r\).

The actual derivative \(Q'\), differentiated at fixed controls,
has the smooth root expression \eqref{eq:Q-outer}. It equals one
at \(r=0\), so \(Q'-1=rU_1(r,m,v,p)\) with \(U_1\) smooth.
The aligned strength cancellation gives \(U_1(0,0,0,p)=0\).
Hence
\begin{equation}\label{sat:outer-slope-bound}
 |Q'-1|\le K r(r+|m|+|v|).
\end{equation}
At \(\zeta_2\), equations~\eqref{sat:curvature-scale} and
\eqref{sat:outer-slope-bound}, together with \(Q=O(r^2)\), yield
\(D'(\zeta_2)=b+O(m^2)\). At the first event \(Q'=1\).
For positive separation that event is zero and \(Q=0\), giving
\(D'(e_1)=b\); for negative separation it is \(\delta\), with
\(Q=\delta\), giving \(D'(e_1)=b+O(|\delta|)\).
The first and last inequalities in \eqref{sat:junction-slopes}
therefore imply
\[
 -Km^2\le sb\le K|\delta|.
\]
The cubic separation bound proves \(|b|\le K|m|^2\).

For the offset and the states, use the event-uniform value bound
\eqref{eq:quartic-global-value} from Lemma~\ref{lem:exact-value-division}.

On the inactive side the exact closing expansion gives
\[
 E(q)=sa+sbq+sc(\lambda)q^2+O(|q|^3).
\]
After shrinking the state interval, its last two terms are bounded
above by \(-\gamma q^2\) for one \(\gamma>0\). Completing the
square gives \(E(q)\le sa+Kb^2\). The positive inactive maximum
from Lemma~\ref{sat:saturation} therefore implies
\begin{equation}\label{sat:offset-lower-coarse}
 sa\ge-Kb^2.
\end{equation}
The next, negative minimum lies between the two actual events, where
\(|q|\le K|\delta|\). Equation~\eqref{eq:quartic-global-value} and the
closing expansion give there
\[
 |D(q)-a|\le K\{ |b\delta|+|\delta|^{3/2}+\delta^2\}.
\]
Consequently
\begin{equation}\label{sat:offset-upper-coarse}
 sa\le K\{ |b\delta|+|\delta|^{3/2}+\delta^2\}
       \le K|m|^{9/2}.
\end{equation}
Together with \eqref{sat:offset-lower-coarse}, this proves
\(|a|\le K|m|^4\).

The three parameter bounds just proved satisfy the hypotheses of
Lemma~\ref{mainloc:tail}, which confines every zero to $|q|<Mm^2$.
Restricting the auxiliary controls to the original observations proves
the assertion for the original family.
\end{proof}

\subsection{The four-root limit}

Fix the sign sector $A_1(0)=B>0$ and $c(0)=-C<0$, and put $h=-b$.
We use the canonical transported strength sum for $m$.
Define
\[
 A_4=C^3a/m^4,\qquad S_4=Ch/m^2,\qquad
 T_4=3BC^2\delta/(2m^3),\qquad x=C\sqrt q/m.
\]
The constants $B,C$ are those of the reference system and are not varied
with the parameters. Write
\[
 \Phi_{A,S,T}(x)=A-Sx^2+x^3-x^4+Tx,
 \qquad
 \mathcal K_4=\{(-e_4(r),e_2(r),e_3(r)):r_i\ge0,\ \sum_i r_i=1\}.
\]

Fix one actual observation inverse $\lambda(a,b,m,\delta;\xi_0)$ on
a neighbourhood $V\subset\mathbb R^4$ of the zero observations, with
the remaining coordinates fixed. Write
\[
 N_\lambda(I)=\#\{q\in I:D_\lambda(q)=0\},\qquad
 v_m(A,S,T)=\left(\frac{m^4A}{C^3},-\frac{m^2S}{C},m,
                         \frac{2m^3T}{3BC^2}\right).
\]
Choose a compact coefficient box $\mathcal B$, independent of $m$,
containing $\mathcal K_4$ in its interior and every normalized local
six-cycle control in $V$. Proposition~\ref{sat:weighted-confinement}
provides such a box after shrinking $V$. The actual main-scale domain is
\begin{equation}\label{def:actual-main-domain}
 \mathscr T_m=\{(A,S,T)\in\mathcal B:v_m(A,S,T)\in V,
                    N_{\lambda(v_m(A,S,T);\xi_0)}(I)=6\}.
\end{equation}
The actual domain, its realization and the local bifurcations use this
inverse. The sequential necessity arguments below allow arbitrary
permitted families and converging background parameters.

For $x>0$ put $F_m(x)=C^3m^{-4}D(m^2x^2/C^2)$. At fixed controls,
\[
 D_q=\frac{m^2}{2Cx}F_m',\qquad
 D_{qq}=\frac{C}{4x^3}(xF_m''-F_m').
\]
Thus the limiting curvature equation is
$x\Phi''-\Phi'=3x^2-8x^3-T=0$, with positive turn $x=1/4$,
turning value $T=1/16$, and nonzero root $x=3/8$ when $T=0$.

\begin{lemma}[Sequential main-scale convergence]\label{lem:k4-sequence}
Let parameters tend to the reference, let $m>0$ tend to zero, and assume
that $(A_4,S_4,T_4)$ stays bounded. If these triples converge to $(A,S,T)$,
then
\[
 \frac{C^3}{m^4}D(m^2x^2/C^2)\longrightarrow\Phi_{A,S,T}(x)
\]
in $C^k$ on every compact subset of $x>0$, for every fixed finite $k$.
No rank condition or prescribed rate of convergence of the background
parameters is required.
\end{lemma}
\begin{proof}
Put $r=mx/C$ and $v=\delta/r^2=2mT_4/(3Bx^2)$. On a positive compact
$x$-interval, $r,m,v\to0$ uniformly with all fixed-order $x$ derivatives.
The exact value division gives
\[
 Q-q=mr^3+\tfrac34m^2r^4+r^5L_0+\delta rH,
 \qquad H(0,0,0,p)=\tfrac32A_1(p,0,0).
\]
Consequently $C^3(Q-q)/m^4\to x^3+Tx$ in every such $C^k$ norm.
In the closing expansion
\[
 D=a+bQ+(Q-q)+c(\lambda)Q^2+Q^3S_R,
\]
we have $Q=q+O_{C^k}(m^4)$, $b=-m^2S_4/C$, and $c(\lambda)\to-C$.
The four remaining terms therefore converge respectively to
$A$, $-Sx^2$, $-x^4$, and zero after multiplication by $C^3/m^4$.
All smooth profiles are evaluated in one common compact neighborhood.
Smooth dependence controls the state derivatives even when the background
parameters converge arbitrarily slowly relative to $m$.
\end{proof}

\begin{lemma}[The event end of a saturated sequence]\label{lem:k4-event}
Let parameters tend to the reference through returns with six fixed points,
and let their normalized coefficients converge to $(A,S,T)$. Denote the five successive critical points of $D$
by $t_1,\ldots,t_5$, and put $x_j=C\sqrt{t_j}/m$ for $j=3,4,5$.
If $T>0$, all three $x_j$ stay in a compact subset of $(0,\infty)$.
If $T=0$, then $x_3\to0$, $A=0$, and every limit of $x_5$ is at least
$3/8$.
\end{lemma}
\begin{proof}
The critical points interlace the six zeros, so their upper bounds follow
from \eqref{sat:weighted-bounds}. The existing offset estimate gives
$a\le K m^{9/2}$, hence $A\le0$.

Suppose first that $T>0$ and that $x_3\to0$ along a subsequence.
At $q=t_3$ both channels are active. Let $y=\Gamma_1(q)-\delta\ge0$.
Since $q^{3/2}/\delta$ is a fixed positive constant times $x_3^3/T_4$,
this ratio tends to zero. Thus
\[
 q-y\ge\delta/2,\quad 0\le y<q,\quad
 \sqrt q-\sqrt y\ge\delta/(4\sqrt q).
\]
The actual derivative expansion, with all coefficients evaluated at the
same physical parameter, is
\[
 Q_q-1=\tfrac32A_1(\sqrt q-\sqrt y)
       +\tfrac32(A_1+A_2)\sqrt y+O(q).
\]
Here $A_1\ge B/2$ and $A_1+A_2=m>0$. Also $Q=O(q)$,
$Q_q=1+O(\sqrt q)$, and $R'(Q)-1=-h+O(q)$. It follows that
\[
 D_q(q)\ge k\delta/\sqrt q-K(h+q).
\]
After multiplying by $C/m^2$, the right-hand side is bounded below by
$k_1T_4/x_3-K_1-K_2x_3^2$, which tends to $+\infty$. This contradicts
$D_q(t_3)=0$. The ordering $x_3<x_4<x_5$ proves the first assertion.

Now suppose $T=0$. Write $r_-<r_+$ for the square roots of the two outer
curvature zeros. They lie on opposite sides of the unique positive
turning point $r_c$, and the actual curvature graph gives
\[
 \begin{aligned}
 \delta&=r^2\{mH_g(r,m,p)+rB_g(r,p)\},\\
 H_g(0,0,p)&=\frac{2}{A_1(p,0,0)},&
 B_g(0,p)&=\frac{16c_0(p)}{3(1+b)A_1(p,0,0)}.
 \end{aligned}
\]
For bounded $x=Cr/m$, this equation becomes
\[
 T_4=3x^2-8x^3+o(1),\qquad Cr_c/m\longrightarrow1/4.
\]
The error is uniform down to $x=0$, because this is an identity between
smooth graph profiles, not an expansion of the singular curvature.
Both $Cr_\pm/m$ are bounded by the critical-point confinement. A limit
of $Cr_-/m$ lies in $[0,1/4]$ and solves $x^2(3-8x)=0$, hence is zero.
A limit of $Cr_+/m$ is at least $1/4$ and solves the same equation, hence
is $3/8$. Since $t_3<r_-^2$ and $t_5>r_+^2$, the location assertions follow.

The event-uniform value estimate yields
\[
 |D(q)-a|\le K\{hq+q^2+mq^{3/2}+m^2q^2+q^{5/2}
                +\delta\sqrt{q+\delta}+\delta q^{3/2}\}.
\]
At $q=t_3$, division by $m^4$ bounds this by a constant times
\[
 x_3^2+x_3^3+x_3^4+m^2x_3^4+mx_3^5
 +T_4x_3+\sqrt m\,T_4^{3/2}+m^2T_4x_3^3,
\]
which tends to zero. Since $D(t_3)>0$, this gives $A\ge0$, and therefore
$A=0$.
\end{proof}

\begin{proposition}[Necessary four-root geometry]\label{prop:k4-necessary}
For every allowed family, each cluster point of the normalized
coefficients of a six-fixed-point sequence tending to the reference
belongs to $\mathcal K_4$.
\end{proposition}
\begin{proof}
Take a convergent coefficient subsequence, using \eqref{sat:weighted-bounds}.
In this sector \eqref{sat:junction-slopes} gives $b=D'(0)<0$ on every
saturated return, so $S_4=-Cb/m^2>0$ and its limit satisfies $S\ge0$.
If $T>0$, Lemmas~\ref{lem:k4-sequence} and~\ref{lem:k4-event} give three
positive limiting critical points $u\le v\le w$ of $\Phi$, with
\[
 \Phi(u)\ge0,\qquad\Phi(v)\le0,\qquad\Phi(w)\ge0.
\]
If $k\le3$ critical points converge to the same positive point $z$,
repeated Rolle gives zeros of $F_m^{(j)}$ tending to $z$ for
$1\le j\le k$. The $C^3$ convergence yields $\Phi^{(j)}(z)=0$ for
those $j$, so $z$ is a zero of $\Phi'$ of multiplicity at least $k$.
If $u<v<w$, the four successive intervals
$[0,u],[u,v],[v,w],[w,\infty)$ each contain a zero, since $\Phi(0)=A\le0$
and $\Phi$ tends to $-\infty$. A zero at a shared critical endpoint has
multiplicity at least two, so the four interval allocations are valid
when counted with multiplicity. If $u=v<w$, then $\Phi(u)=0$ and the
double derivative zero gives a triple polynomial zero at $u$; a fourth
zero lies in $[w,\infty)$. The case $u<v=w$ is analogous. If all three
merge, the common value is zero and the root has multiplicity four.
Thus all four roots are nonnegative and their sum is one.

If $T=0$, the event lemma gives $A=0$. After subselection $x_5\to w\ge3/8$,
and positive-compact convergence gives $\Phi(w)\ge0$. Therefore
$-S+w-w^2\ge0$, so $0\le S\le1/4$. The factorization
$\Phi=x^2(-x^2+x-S)$ again gives four nonnegative roots with sum one.
\end{proof}

\begin{proof}[Completion of Theorem~\ref{thm:main-saturation-summary}]
The interior of $\mathcal K_4$ consists exactly of coefficient triples
with four distinct positive roots. Such triples are open by simple-root
continuation. A zero root or a multiple root is a boundary point: a
small constant-term perturbation gives respectively a negative root or
violates one of the necessary critical-value signs for four real roots.
Triples with distinct positive roots are dense, by perturbing the root
vector inside the simplex while preserving its sum.

For four distinct positive roots $r_i$, the positive limiting return is
$-\prod_i(x-r_i)$. On the negative scaled state $y=C^2q/m^2$, the limiting
return is $-(y^2+e_2y+e_4)$. The six pair products satisfy
$e_2\ge6\sqrt{e_4}$ by AM--GM; hence this quadratic has two distinct
negative roots. On a compact subset of the interior, all six roots have
uniform isolation and nonzero derivative margins. The existing mixed
$C^1$ weighted convergence and the fixed observation inverse continue
these roots to actual simple fixed points, uniformly in the compact
coefficient set. The common six-root upper bound excludes additional
fixed points in the physical interval.

For the Hausdorff upper inclusion, a contrary sequence of closure points
can be approximated within $1/n$ by true points of $\mathscr T_{m_n}$.
Compactness of $\mathcal B$ and Proposition~\ref{prop:k4-necessary} then
force its limit into $\mathcal K_4$. For the lower inclusion, approximate
compact $\mathcal K_4$ by a finite net of interior points and realize all
of them at every sufficiently small $m$. The fixed coarse box is chosen
using \eqref{sat:weighted-bounds} and enlarged to contain $\mathcal K_4$ in its
interior. The box also includes the normalized controls of any saturated
sequence with small positive offset; \eqref{sat:offset-upper-coarse}
forces their offset coordinate to tend to zero.
\end{proof}

\subsection{Boundary geometry of the limiting domain}

Figure~\ref{fig:main-saturation-sections} displays two sections of the
limiting coefficient region; the critical-value criterion below describes
their boundaries.

\begin{figure}[!htbp]
\centering
\includegraphics[width=\textwidth]{figures/main-saturation-sections.pdf}
\caption{Sections of the limiting four-root domain at $S_4=0.2$ and
$S_4=0.35$, with separate axis ranges and the displayed ordinate scaling.
The boundaries are drawn from the critical-value interval described below. In the
left panel the segment $A_4=0$ represents a zero root. The remaining boundary
arcs have repeated roots; the change of lower arc occurs when the two local
maxima have equal value. These plots illustrate the algebraic limiting set,
not the corrected bifurcation boundary of a nonzero-scale actual return.}
\label{fig:main-saturation-sections}
\end{figure}

The symmetric mean inequalities give
\[
 -\frac1{256}\le A_4\le0,\qquad
 0\le S_4\le\frac38,\qquad 0\le T_4\le\frac1{16}.
\]
All three nonzero extrema occur at $r_i=1/4$, where
$\Phi=-(x-1/4)^4$. At the opposite face, $e_3=0$ forces at least two
roots to vanish, and therefore
\[
 \mathcal K_4\cap\{T_4=0\}
 =\{(0,S_4,0):0\le S_4\le1/4\}.
\]
These are exact bounds for the limiting set, and hence asymptotic bounds
along saturated sequences. The finite-scale domain retains its actual
boundary corrections.

For fixed $S_4,T_4$ with three positive critical points $u<v<w$, put
$f(x)=-S_4x^2+x^3-x^4+T_4x$. The admissible offset interval is
\[
 \max\{-f(u),-f(w)\}\le A_4\le\min\{0,-f(v)\}.
\]
It is nonempty precisely when these critical-value inequalities are
compatible. Their closures include merged critical points.
